\section{Synthesis: Jurisdiction, Document Structure, and Multilingual Legal AI}
\label{sec:jurisdiction}

This chapter draws together a theme this survey has touched at several points along its chronology,
most directly through \citet{ref312025} and Legalyze \citep{hamza2026} in Section~\ref{sec:2026}:
legal information retrieval is jurisdiction-dependent in a way general-purpose retrieval is not. The
relevance of a document depends not only on its topic but on its legal authority, jurisdiction,
temporal validity, and relationship to the question asked. A question on Pakistani criminal
procedure should prioritize applicable Pakistani legislation and judicial decisions. A semantically
similar provision from another jurisdiction may help explain a general concept, but a system should
not automatically treat it as governing authority. Jurisdiction-sensitive retrieval therefore has to
consider territorial and institutional jurisdiction, authority type (legislation, binding judgment,
persuasive judgment, commentary), temporal validity, procedural context, and precedential
relationships. An embedding model, of the kind described in Section~\ref{sec:rag2024}, can mark two
passages as highly similar without determining whether either is legally applicable. These
constraints should shape retrieval and verification directly, not just appear as explanatory text in
the final answer.

Legal documents also carry multiple layers of organization. Statutes divide into parts, sections,
subsections, and provisos. Judgments contain procedural history, issues, reasoning, and final
orders. Fixed-size chunking can separate a provision from its heading, definition, or exception,
producing a fragment that remains readable while losing the context legal interpretation needs. The
structure-aware retrieval techniques described in Section~\ref{sec:rag2024} address this in four
ways. Hierarchical metadata links each passage to its document, chapter, and section.
Context-preserving chunking keeps headings and exceptions attached. Parent-child retrieval recovers
a passage's broader context on demand. Provenance preservation records the source document and
version a reader needs to verify it. Pakistan shows why this matters in practice. Its legal
information needs span the Constitution, the Pakistan Penal Code, the Code of Criminal Procedure,
the Qanun-e-Shahadat Order, and judicial decisions that serve different legal functions, so a corpus
can be large and still miss relevant amendments or authoritative judgments. A provision's wording or
interpretation can also change over time. A system has to track this by recording source dates and
applicable versions, instead of assuming the newest document is the most authoritative.
\citet{ref312025} demonstrated a RAG-based framework combining semantic and keyword retrieval over
selected Pakistani legal materials. Legalyze \citep{hamza2026} extended this with hierarchical
indexing, hybrid retrieval, and rank fusion, as Section~\ref{sec:2026} described. A meaningful
comparison between the two would still need a shared evaluation dataset, consistent query sets, and
common retrieval and answer-quality measures.

A further challenge arises when a user's query language differs from the language or script of the
authoritative source. This is common when Pakistani legal questions appear in Urdu or Roman Urdu
while the relevant statute or judgment appears in English. Direct translation alone can fall short,
because legal concepts carry jurisdiction-specific meanings that do not map exactly across
languages. Table~\ref{tab:multilingual} summarizes the main technical approaches.

\begin{table*}[t]
\centering
\caption{Technical approaches to multilingual legal retrieval.}
\label{tab:multilingual}
\footnotesize
\setlength{\tabcolsep}{5pt}
\renewcommand{\arraystretch}{1.25}
\begin{tabularx}{\textwidth}{>{\raggedright\arraybackslash}p{3.0cm} >{\raggedright\arraybackslash}X >{\raggedright\arraybackslash}X}
\toprule
\textbf{Approach} & \textbf{Main advantage} & \textbf{Main limitation} \\
\midrule
Query translation & Enables use of existing monolingual retrieval infrastructure. & Translation may distort legal terms or procedural meaning. \\
Multilingual dense retrieval & Supports cross-language semantic matching in a shared embedding space. & Similar embeddings do not guarantee legal-concept equivalence. \\
Transliteration normalization & Helps retrieve Roman Urdu and Urdu spelling variants. & Many transliterations remain ambiguous or inconsistent. \\
Legal terminology alignment & Links expressions to jurisdiction-specific legal concepts. & Requires curated terminology and entity-linking resources. \\
\bottomrule
\end{tabularx}
\end{table*}

A practical architecture combines these methods instead of relying on just one. Language
identification and transliteration normalization process the query first. Legal entity linking and
terminology alignment then identify the relevant statutes and concepts. Retrieval searches both the
original and normalized query before rank fusion and jurisdictional filtering. The system should
keep the original query and the original authoritative passage throughout. Translation can support
retrieval, but the final answer should not quietly substitute a translated interpretation for the
governing source text. We should also carry out evaluation separately for each supported language,
rather than assume it carries over from an English-only benchmark. Jurisdiction, document structure,
and language are consequently tied together, not separable concerns. Language normalization can help
locate a relevant provision, but document structure still has to preserve its context. Hierarchical
indexing can identify the correct section, but jurisdictional and temporal checks remain necessary
to work out whether it actually applies. For Pakistani legal AI, the resulting requirement is clear:
effective retrieval needs more than semantic similarity. It needs legal hierarchy, authority
relationships, temporal metadata, source provenance, and cross-language concept alignment, all
working together. That requirement leads into the open research directions we discuss next.
